% ==============================================================================
% abstract.tex
% ==============================================================================

\begin{abstract}
Reconstructing three-dimensional computed tomography (CT) volumes from a small number of X-ray projections is a critical problem in medical imaging, offering potential reductions in radiation dose, scan time, and hardware cost. This survey provides a systematic review of learning-based methods for few-view X-ray-to-CT reconstruction, covering the period from 2018 to 2025. We propose a taxonomy that organises existing approaches into four major families: (1)~generative adversarial network (GAN)-based methods, (2)~transformer-based methods, (3)~neural field and implicit representation methods, and (4)~diffusion-based generative methods. For each family, we analyse representative architectures, training strategies, and reported performance. We further review commonly used datasets, discuss the persistent domain gap between synthetic and clinical data, and present a comparative analysis of evaluation metrics. Finally, we identify open challenges and outline promising directions for future research, including scalable diffusion models, clinically validated benchmarks, and hybrid reconstruction pipelines.
\end{abstract}
